\Needspace{13\baselineskip}
\section{R2: Carrier invention and accumulator discovery}\label{sec:carriers}

\Q{R2.Q1}{What finite examples can decide about a minimal carrier}{Is there a constructive procedure that finds a smallest carrier type from a grammar, or proves that no such carrier fits the examples?}

\answer There is a complete constructive answer for an explicitly finite problem. There is no corresponding general decision procedure merely because the \emph{type grammar} is finite. The carrier, step, initialization, and observation map must be distinguished:
\begin{equation}\label{eq:fold-factor}
 q([])=z,\qquad q(a::x)=s(a,q(x)),\qquad h=o\circ q,
\end{equation}
where $z:K$, $s:A\to K\to K$, and $o:K\to T$. Minimal state cardinality, smallest syntax for $K$, cheapest executable realization, and best match to an intended function are four different objectives.

If $K$ ranges over a finite list of types but step bodies range over arbitrary Lean terms, type and contract synthesis have not become a finite problem. If $o$ is unrestricted and $K=\List A$ is allowed, every function $h$ factors through the identity list fold: choose $z=[]$, $s=\mathsf{cons}$, and $o=h$. Consequently, a semantic obstruction to writing $h$ itself as a fold cannot automatically rule out a fold followed by an arbitrary finisher.

Here is a precise finite version. Fix finite alphabets $A$, finite output sets $T$, explicitly enumerated finite carriers $K$, and permit every total table for $s$ and $o$. For each $K$, enumerate its initial state, transition table, and output table; evaluate all sample words. Order carriers by a declared cost and return the first satisfying representation. This terminates, and failure is a genuine finite-state impossibility result for these assumptions. Alternatively encode the same problem as a finite constraint satisfaction problem, with variables for the state of each observed suffix and the entries of the two tables. A complete Lean-native backtracking solver suffices; an external solver is optional, not necessary.

To translate this answer back to syntax, require either an emitter that realizes every table using the carrier's finite enumeration and equality, or an explicit finite grammar for the initialization, step, and finisher. The latter can be exhaustively searched by total syntactic size. A carrier table that is extensionally realizable but too large for the step grammar does not solve the bounded syntactic problem. For arbitrary Lean checks, resource exhaustion must remain unknown; a decidable finite fragment or a completed certificate checker is needed for a definitive negative answer.

\subsubsection*{The exact semantic lower bound}

The following formulation makes the automaton analogy precise without confusing samples with a complete specification. It is a direct algebraic argument, related to classical fold characterization and automata learning.\cite{folds,angluin}

\begin{definition}[Continuation equivalence]
For $h:\List A\to T$, define
\[
 x\sim_h y\quad\Longleftrightarrow\quad
 \forall u:\List A,\ h(u\mathbin{+\!+}x)=h(u\mathbin{+\!+}y).
\]
These are \emph{left} continuations because $q$ in \cref{eq:fold-factor} is a right fold.
\end{definition}

\begin{theorem}[Minimal reachable state quotient]\label{thm:carrier}
Every factorization \cref{eq:fold-factor} has at least as many reachable states as $\sim_h$ has equivalence classes. Conversely, the quotient by $\sim_h$ supports a factorization of $h$. Thus it has minimum reachable-state cardinality among arbitrary set-theoretic right-fold representations with an output map.
\end{theorem}
\begin{proof}
If $q(x)=q(y)$, induction on $u$ gives $q(u\mathbin{+\!+}x)=q(u\mathbin{+\!+}y)$. Applying $o$ shows $x\sim_h y$. Hence the map from a reachable state $q(x)$ to the equivalence class $[x]$ is well defined and surjective.

For the converse, use states $[x]$, initial state $[\,]$, transition $s(a,[x])=[a::x]$, and output $o([x])=h(x)$. Output is well defined by choosing $u=[]$. The transition is well defined because $u\mathbin{+\!+}(a::x)=(u\mathbin{+\!+}[a])\mathbin{+\!+}x$. Induction gives $q(x)=[x]$, so $o(q(x))=h(x)$.
\end{proof}

For infinite cardinal comparisons, the theorem is stated in ordinary classical set theory with choice; its finite lower-bound applications need only elementary finite counting. It constructs a mathematical quotient, not a computable presentation of that quotient from finitely many examples. In particular, it does not identify a preferred Lean type among isomorphic carriers. Infinite types such as \code{Nat} also make raw cardinality a poor proxy for useful state structure: encodings can hide considerable computation in the transition and finisher.

Finite observations do nevertheless yield certified lower bounds. Choose suffixes $x_i$ and prefixes $u_j$ for which outputs are specified. Join two suffixes by an edge when some shared observed prefix gives different outputs. Any consistent carrier assignment must color this incompatibility graph: adjacent suffixes need distinct states. A witnessed clique of size $r$ proves that fewer than $r$ states cannot work. The chromatic number is a stronger lower bound; transition consistency can strengthen it further.

Do not count distinct \emph{partial} observation rows as distinct states. Compatibility of incomplete rows is not transitive. For example,
\[
 (0,?),\qquad (?,1),\qquad (1,1)
\]
has compatible first--second and second--third pairs, but incompatible first--third entries. Unknown cells are not a new output value. Nor does enumerating the values constructible up to depth $d$ provide an upper bound on all values a type can hold: it gives a lower bound on its discovered inhabitants. Cardinality pruning requires an actual upper-bound proof.

\begin{example}[A four-state requirement from fifteen observations]\label{ex:fourstate}
Let $A=\{a,b\}$, and let $h(x)$ say that $x$ contains both letters. The suffixes $[],[a],[b],[a,b]$, observed under prefixes $[],[a],[b]$, give the rows
\[
\begin{array}{c|ccc}
 x & [] & [a] & [b]\\\hline
 {[]}&0&0&0\\
 {[a]}&0&0&1\\
 {[b]}&0&1&0\\
 {[a,b]}&1&1&1
\end{array}
\]
so all four suffixes are pairwise incompatible. All observations involved are words of length at most three. No carrier with at most three states fits them. A pair of Booleans works: record whether $a$ has appeared and whether $b$ has appeared; combine the two bits for the output. The supplied script independently exhausts every one-, two-, and three-state transition/output table with a fixed initial-state label and finds no survivor.
\end{example}

For Leant2 this suggests a dual output from carrier inference: a proposed carrier plus code, and, when available, a small distinguishing-context certificate excluding smaller finite carriers. That is more informative than a heuristic carrier guess and much cheaper to check than rerunning the whole search.

\Q{R2.Q2}{Polymorphic automata and canonical inputs}{Does the state-machine perspective extend to polymorphic functions, perhaps using positions or equality patterns as the alphabet?}

\answer Yes, under an explicit relational-parametricity assumption for the program fragment and its providers. It is not valid for every term allowed by Lean's standard axiom profile. The literature on testing polymorphic properties establishes reductions to suitable monomorphic instances for specified classes of parametrically polymorphic properties; its hypotheses are substantive.\cite{polytests}

Suppose $f_A:\List A\to\List A$ is natural in $A$:
\begin{equation}\label{eq:natural}
 f_B(\mathsf{map}\ e\ x)=\mathsf{map}\ e\,(f_A(x))
 \quad(e:A\to B).
\end{equation}
Let $c_n=[0,\ldots,n-1]:\List(\mathsf{Fin}\ n)$. For any list $x$ of length $n$, let $e_x:\mathsf{Fin}\ n\to A$ select its entries. Then
\[
 f_A(x)=\mathsf{map}\ e_x\,(f_{\mathsf{Fin}\ n}(c_n)).
\]
Thus a list of output positions describes every behavior at a fixed input length, including duplication, deletion, and reordering. The case $n=0$ is included: the finite type has no inhabitants, so its output list is empty. Agreement on canonical inputs through length $N$, together with \cref{eq:natural} for both functions, proves agreement on all types and all lists of length at most $N$. It does not prove agreement at greater lengths.

For contrast, in classical Lean one can define a noncomputable family that returns its input when $A$ is a subsingleton and returns the empty list otherwise, using classical decidability of \code{Subsingleton A}. Mapping from a two-element type to a singleton shows failure of \cref{eq:natural}. The type $\forall A,\List A\to\List A$ alone therefore does not license the canonical-position argument.

If the language exposes lawful decidable equality on elements, arbitrary maps are no longer the relevant relation-preserving morphisms. Equality patterns matter: canonical inputs should realize partitions of positions, not only the all-distinct partition. Extra operations, orders, algebraic dictionaries, or distinguished elements require richer test structures and preservation laws. A user-supplied \code{BEq} without laws cannot be silently identified with propositional equality.

The practical first implementation should certify parametricity for a deliberately small grammar of constructors, eliminators, and explicitly approved polymorphic providers. Until that exists, positional probes remain useful heuristics for presentation and counterexample discovery, but must not justify permanent merging or impossibility claims. Symbolic state representations may keep a finite control tag together with stored positions or elements; a finite abstract alphabet does not imply that the concrete carrier itself is finite.

\Q{R2.Q3}{Backward fusion and auxiliary state}{Can examples be worked backward to discover a carrier, and is the method complete for useful classes?}

\answer Use backward reasoning to propose a \emph{commuting algebra}, not to infer a unique hidden auxiliary function from samples. Write an unknown summary $q$ and seek equations
\[
 q([])=z,\quad q(a::x)=s(a,q(x)),\quad o(q(x))=h(x).
\]
When $q$ is a tuple $(q_1,\ldots,q_m)$, failed closure obligations often reveal which additional summary component is missing. The state grammar and the equations for all components should be searched together. Auxiliary state is useful precisely when it makes the local transition inexpensive and expressible.

For natural-number Fibonacci, the summary
\[
 q(n)=(F_n,F_{n+1}),\quad q(0)=(0,1),\quad s(x,y)=(y,x+y)
\]
turns two-step history into a one-step transition. Its invariant is proved by induction from the Fibonacci equations. For maxima of possibly empty lists, an optional stored element avoids inventing an element of an arbitrary empty type. For a tail-recursive helper, a strengthened equation such as $\mathsf{go}(x,a)=h(x)\star a$ simultaneously identifies the accumulator and the induction invariant.

Reversal clarifies what the optimization does and does not establish:
\[
 \mathsf{reverse}=\fold\ (\lambda a\ r.\ r\mathbin{+\!+}[a])\ [].
\]
This is already a fold with carrier $\List A$. Appending to the end repeatedly gives quadratic list traversal in the usual immutable singly linked representation. A function carrier admits
\[
 q([])=\id,\qquad q(a::x)=\lambda y.\ q(x)(a::y),\qquad
 o(k)=k([]).
\]
The invariant $q(x)(y)=\mathsf{reverse}(x)\mathbin{+\!+}y$ proves correctness. Under constant-time cons and closure construction, it avoids repeated append traversal. Alternatively, changing to a left-fold scheme also permits a linear reversal with a list carrier. Neither observation proves that reversal semantically requires a higher-order carrier.

For a fixed finite alphabet and finite carrier, the exhaustive procedure in R2.Q1 is complete for all finite transition tables, including those suggested by backward reasoning. For a finitely specified polynomial datatype, fixed carrier, and bounded algebra grammar, constructor-by-constructor enumeration is complete relative to that grammar. Once carrier types, helper functions, semantic invariants, and unrestricted finishers can all be invented, no comparable terminating general guarantee follows. A useful promise is fair enumeration by a combined size bound, with backward fusion changing the order and providing pruning certificates when available.

\Q{R2.Q4}{Enumerating recursion motives with parameters}{Can primitive-recursion motives be enumerated completely and efficiently without recreating the natural-number search explosion?}

\answer Completeness and efficiency need separate statements. For a finite grammar of motives, local parameters, and auxiliary carriers, enumerate by a \emph{joint} grade rather than running an unbounded motive search before body synthesis. Use a cost such as
\[
 |K|+|M|+|z|+|s|+|o|+\lambda\,\#\text{recursive schemes}.
\]
The notation denotes syntax size, not state cardinality. Fairly increasing this grade eventually considers every term in the grammar; it does not bound the time needed at a given grade.

Start with the result family, function carriers over existing parameters, small products, and options. For an indexed datatype, abstract the indices in the result family and generalize dependency-closed context subsets, as detailed in R4. For ordinary \code{Nat.rec}, try outermost recursion first and charge nested recursion explicitly. The uncontrolled multiplication of recursion choices in every natural-number hypothesis is a search-order problem, not an argument for permanently removing all natural-number recursion.

A bound on carrier syntax alone is insufficient because a function carrier can conceal arbitrarily large step code. Conversely, limiting recursion to an outermost scheme is an effective first tier but is incomplete for the unrestricted grammar. Make the tier visible in the completed-bound metadata and eventually expand it in the reference lane. No theorem turns arbitrary higher-order motive invention into a cheap finite preprocessing step.

\Q{R2.Q5}{Using inductive generalization and lemma discovery}{What can techniques from inductive theorem proving contribute to carrier and helper invention?}

\answer Transfer the \emph{shape of the obligation} as much as the final theorem-proving machinery. If an induction step is blocked because an accumulator has changed, generalize that accumulator before induction. If the step repeatedly creates an unexpressed statistic, propose that statistic as a component of the carrier. If normalization exposes an associative combination, propose a helper invariant that distributes over it.

Theory exploration systems such as Hipster already combine conjecture generation with induction-based proof attempts, including discovery of lemmas missing from a current proof.\cite{hipster} The Leant2 adaptation should make this a bounded service: extract terms and type-correct anti-unifications from the failed preservation goal; propose a small number of generalized equations; refute bad proposals using concrete instances; then prove surviving equations before using them as rewrite rules or search cuts. A theorem that is useful only after a helper is introduced should be linked to that helper's synthesis task, not asserted as an oracle.

Rippling-style annotations can be represented as search metadata describing how a goal differs from its induction hypothesis. They suggest which term position to abstract or which lemma might move the difference outward. The resulting suggestion is not a theorem and should not force a single irreversible generalization. In a native Lean implementation, the first version can be much smaller than a full rippling engine: recognize accumulator arguments, retained input tails, and constructor-preservation equations, then try a bounded library of invariant schemas. Measure the incremental solved tasks and the extra proof debt, not merely the number of conjectures generated.
